Meeting stringent Time-To-First-Token (TTFT) requirements is crucial for LLM applications. To improve efficiency, modern LLM serving systems adopt disaggregated architectures with diverse parallelisms, introducing complex multi-stage workflows involving reusable KV-block retrieval, collective communication, and P2D transfer. Flows from dependent stages overlap within and across requests on shared bottleneck links, making TTFT highly susceptible to network contention and necessitating stage-aware scheduling. Unfortunately, most existing works schedule flows in a stage-agnostic manner, leading to uncoordinated contention that constitutes a primary cause of SLO violations.

In this paper, we present \sys, a holistic \underline{m}ulti-stage \underline{f}low \underline{s}cheduling mechanism designed to maximize TTFT SLO attainment. At its core, \sys approximates the Least-Laxity-First (LLF) scheduling policy without requiring precise knowledge of a request's remaining slack. It achieves this through a \emph{Defer-and-Promote} principle implemented through a Reverse Multi-Level Queue (RMLQ) structure. By dynamically promoting task precedence as effective laxity diminishes, \sys prioritizes flows with less laxity while preventing requests with loose SLOs from prematurely consuming network bandwidth. We implement \sys as a pluggable module integrated into vLLM, and evaluate it on a 8-server, 32-GPU testbed as well as through large-scale simulations. Our results demonstrate that \sys effectively outperforms state-of-the-art baselines, improving the TTFT SLO attainment by 1.2$\times$--2.4$\times$.
